\subsection{p-adic Lie groups.}

Such groups are met for the first time in 1907 in the work of Hensel {[}157 e{]} on the p-adic analytic functions (defined by expansions in entire series). He studies notably the exponential and the logarithm; in spite of the a priori surprising behaviour of the series that define them (for example the exponential series does not converge everywhere), their functional properties remain valid, which supplies a local isomorphism between the additive group and the multiplicative group of \(Q_{p}\) (or, more generally, of every complete ultrametric field of characteristic zero).

It is equally about commutative groups (but not linear this time) that the work of A. Weil {[}330 a{]} and E. Lutz {[}210{]} on the elliptic p-adic curves (1936) is concerned. Other than the arithmetic applications, there can be found the construction of a local isomorphism of the group with the additive group, based on the integration of an invariant differential form. This method can be equally applied to abelian varieties, as is remarked shortly afterwards by C. Chabauty who uses it without any more explanation in order to prove a particular case of the ``Mordell conjecture'' {[}61{]}.

From this moment, it was clear that the local theory of Lie groups can be applied almost without change to the p-adic case. The fundamental theorems of the ``dictionary'' Lie groups-Lie algebras are established in 1942 in the thesis of R. Hooke {[}166{]}, a pupil of Chevalley; this work contains also the p-adic analogue of the theorem of E. Cartan on the closed subgroups of real Lie groups.

More recently, M. Lazard {[}195 b{]} develops a more precise form of the ``dictionary'' for compact analytic groups over \(\mathbf{Q}_p\) . He shows that the existence of an analytic \(p\) -adic structure on a compact group \(G\) is tightly linked to that of certain filtrations on \(G\) , and gives various applications of it (for example to the cohomology of \(G\) ). One of Lazard's tools is an improvement of Dynkin's results on the convergence of the \(p\) -adic Hausdorff series {[}195 a{]}.

